\documentclass[12pt]{article}
\usepackage[utf8]{inputenc}
\usepackage[T1]{fontenc}
\usepackage{hyperref}
\title{Vemgo-Mabas Language Resources for K-12 Curriculum Development}
\author{Compiled by Claude Sonnet 5 (claude-sonnet-5) for Andr\'as Kornai}
\date{}
\begin{document}
\maketitle
\section{Introduction}
Vemgo-Mabas is a Chadic (Biu-Mandara) language of northeastern Nigeria (Madagali area) and northern Cameroon. Speaker estimates vary widely (about 11,000 on Wikipedia; 1,000 to 1,700 for Cameroon in other sources). Vemgo and Mabas may be separate languages according to Blench. No educational materials were found.
\section{Primary Educational Resources}
None found.
\section{Linguistic and Reference Materials}
\begin{itemize}
\item Wikipedia, Vemgo-Mabas language: \url{https://en.wikipedia.org/wiki/Vemgo-Mabas_language}
\item Mandara Mountains site, The Mabas, Vizik \& Vemgo: \url{https://mandaras.info/MabasVizikVemgo.html}
\item SIL Cameroon, sociolinguistic survey of Glavda, Cineni and Vemgo-Mabas (rapid appraisal): \url{https://www.silcam.org/resources/archives/47435}
\item Numerals comparison page: \url{https://lingweb.eva.mpg.de/channumerals/Vemgo-Mabas.htm}
\item ASJP word list entry: \url{https://asjp.clld.org/languages/VEMGO_MABAS_2}
\end{itemize}
\section{Religious and Community Resources}
\begin{itemize}
\item Joshua Project, Vemgo-Mabas in Cameroon: \url{https://www.joshuaproject.net/index.php/people_groups/13173/CM}
\item Joshua Project, Vemgo-Mabas, Visik in Nigeria: \url{https://joshuaproject.net/index.php/people_groups/13173/NI}
\end{itemize}
\section{Digital and Audio Resources}
\begin{itemize}
\item Global Recordings Network entry: \url{https://globalrecordings.net/en/language/vem}
\end{itemize}
\section{Curriculum Development Context}
Little exists. A curriculum would need basic documentation and orthography work first. The sociolinguistic survey is the best starting point for community context.
\end{document}